\ifdefined\gkver\else\def\gkver{student}\fi
\documentclass[\gkver]{gaokaozhenti}
\begin{document}

\chapter{2008年全国理综Ⅰ}
%% paperType: 理综物理部分

\begin{enumerate}
\item
%% number: 14
%% paperName: 2008年全国理综Ⅰ第 14 题 6 分
%% typeId: 1
%% type: 单选题
%% chapter: 曲线运动
%% point: 平抛运动
%% method: 无
%% score: 6
%% degree: 650
%% duplicateId: 0
%% body:
如图所示，一物体自倾角为 $\theta$ 的固定斜面顶端沿水平方向抛出后落在斜面上。物体与斜面接触时速度与水平方向的夹角 $\varphi$ 满足\xzanswer{D}

\onepicture[w=6.0cm]{figs/14.png}

\fourchoices[answer=D]
{$\tan\varphi=\sin\theta$}
{$\tan\varphi=\cos\theta$}
{$\tan\varphi=\tan\theta$}
{$\tan\varphi=2\tan\theta$}

%% answer: D
%% memo:
\memoanswer{
竖直速度与水平速度之比为 $\tan\varphi=\frac{gt}{v_0}$，竖直位移与水平位移之比为 $\tan\theta=\frac{0.5gt^2}{v_0t}$，故 $\tan\varphi=2\tan\theta$，D 正确。
}

\item
%% number: 15
%% paperName: 2008年全国理综Ⅰ第 15 题 6 分
%% typeId: 2
%% type: 多选题
%% chapter: 运动和力
%% point: 牛顿第二定律
%% method: 无
%% score: 6
%% degree: 650
%% duplicateId: 0
%% body:
如图，一辆有动力驱动的小车上有一水平放置的弹簧，其左端固定在小车上，右端与一小球相连，设在某一段时间内小球与小车相对静止且弹簧处于压缩状态，若忽略小球与小车间的摩擦力，则在此段时间内小车可能是\xzanswer{AD}

\onepicture[w=5.5cm]{figs/15.png}

\fourchoices[answer=AD]
{向右做加速运动}
{向右做减速运动}
{向左做加速运动}
{向左做减速运动}

%% answer: AD
%% memo:
\memoanswer{
对小球水平方向受到向右的弹簧弹力 $N$，由牛顿第二定律可知，小球必定具有向右的加速度，小球与小车相对静止，故小车可能向右加速运动或向左减速运动。
}

\item
%% number: 16
%% paperName: 2008年全国理综Ⅰ第 16 题 6 分
%% typeId: 1
%% type: 单选题
%% chapter: 机械波
%% point: 波的图象
%% method: 无
%% score: 6
%% degree: 650
%% duplicateId: 0
%% body:
一列简谐横波沿 $x$ 轴传播，周期为 $T$，$t=0$ 时刻的波形如图所示。此时平衡位置位于 $x=3\Um$ 处的质点正在向上运动，若 a、b 两质点平衡位置的坐标分别为 $x_{a}=2.5\Um$，$x_{b}=5.5\Um$，则\xzanswer{C}

\onepicture[w=5.5cm]{figs/16.png}

\fourchoices[answer=C]
{当 a 质点处在波峰时，b 质点恰在波谷}
{$t=T/4$ 时，a 质点正在向 $y$ 轴负方向运动}
{$t=3T/4$ 时，b 质点正在向 $y$ 轴负方向运动}
{在某一时刻，a、b 两质点的位移和速度可能相同}

%% answer: C
%% memo:
\memoanswer{
由图可看出波长为 $4\Um$，$t=0$ 时刻 $x=3\Um$ 处的质点向上振动，可得该波向左传播。将整个波形图向左平移 $1.5\Um$ 时，a 质点到达波峰，此时 b 质点正好在平衡位置，与 $t=0$ 时刻平衡位置在 $7\Um$ 处的质点振动状态一样，故 a 质点到达波峰时，b 质点正在平衡位置并向上振动，A 错；将图像整体向左平移 $1\Um$，即波传播 $T/4$ 时，a 的振动状态与 $t=0$ 时刻平衡位置在 $3.5\Um$ 处的质点振动状态一样，即处在平衡位置上方并向 $y$ 轴正方向运动，B 错；将图像整体向左平移 $3\Um$，即波传播 $3T/4$ 时，a 的振动状态与 $t=0$ 时刻平衡位置在 $9.5\Um$ 处和 $1.5\Um$ 的质点振动状态一样，即处在平衡位置下方并向 $y$ 轴负方向运动，C 对；a、b 质点相隔 $3\Um$，即相差 $3T/4$，速度相同的质点应该在半周期内才会出现，故 D 错。
}

\item
%% number: 17
%% paperName: 2008年全国理综Ⅰ第 17 题 6 分
%% typeId: 1
%% type: 单选题
%% chapter: 曲线运动
%% point: 天体运动
%% method: 无
%% score: 6
%% degree: 450
%% duplicateId: 0
%% body:
已知太阳到地球与地球到月球的距离的比值约为 390，月球绕地球旋转的周期约为 27 天。利用上述数据以及日常的天文知识，可估算出太阳对月球与地球对月球的万有引力的比值约为\xzanswer{B}

\fourchoices[answer=B]
{0.2}
{2}
{20}
{200}

%% answer: B
%% memo:
\memoanswer{
设太阳质量 $M$，地球质量 $m$，月球质量 $m_{0}$，日地间距离为 $R$，月地间距离为 $r$，日月之间距离近似等于 $R$，地球绕太阳的周期 $T$ 约为 360 天，月球绕地球的周期为 $t=27$ 天。对地球绕着太阳转动，由万有引力定律 $G\frac{Mm}{R^2}=m\frac{4\pi^2R}{T^2}$，同理对月球绕着地球转动 $G\frac{mm_0}{r^2}=m_0\frac{4\pi^2r}{t^2}$，则太阳质量与地球质量之比为 $M:m=\frac{R^3T^2}{r^3t^2}$；太阳对月球的万有引力 $F=G\frac{Mm_0}{R^2}$，地球对月球的万有引力 $f=G\frac{mm_0}{r^2}$，故 $F:f=\frac{Mr^2}{mR^2}$，带入太阳与地球质量比，计算出比值约为 2，B 对。
}

\item
%% number: 18
%% paperName: 2008年全国理综Ⅰ第 18 题 6 分
%% typeId: 2
%% type: 多选题
%% chapter: 原子和原子核
%% point: 人工核反应
%% method: 无
%% score: 6
%% degree: 650
%% duplicateId: 0
%% body:
三个原子核 X、Y、Z，X 核放出一个正电子后变为 Y 核，Y 核与质子发生核反应后生成 Z 核并放出一个氦核（$\ce{^4_2 He}$），则下面说法正确的是\xzanswer{CD}

\fourchoices[answer=CD]
{X 核比 Z 核多一个质子}
{X 核比 Z 核少一个中子}
{X 核的质量数比 Z 核质量数大 3}
{X 核与 Z 核的总电荷是 Y 核电荷的 2 倍}

%% answer: CD
%% memo:
\memoanswer{
设原子核 X 的质量数为 $x$，电荷数为 $y$，依题意写出核反应方程，根据质量数守恒和电荷数守恒，可得原子核 Y 的质量数为 $x$、电荷数为 $y-1$，原子核 Z 的质量数为 $x-3$、电荷数为 $y-2$。由此可得 X 核的质子（$y$）比 Z 核的质子（$y-2$）多 2 个，A 错；X 核的中子（$x-y$）比 Z 核的中子（$x-y-1$）多 1 个，B 错；X 核的质量数（$x$）比 Z 核的质量数（$x-3$）多 3 个，C 对；X 核与 Z 核的总电荷（$2y-2$）是 Y 核电荷（$y-1$）的 2 倍，D 对。
}

\item
%% number: 19
%% paperName: 2008年全国理综Ⅰ第 19 题 6 分
%% typeId: 1
%% type: 单选题
%% chapter: 气体性质
%% point: 阿伏伽德罗常数
%% method: 近似与估算
%% score: 6
%% degree: 450
%% duplicateId: 0
%% body:
已知地球半径约为 $6.4\times10^{6}\Um$，空气的摩尔质量约为 $29\times10^{-3}\Ukgmol$，一个标准大气压约为 $1.0\times10^{5}\UPa$，利用以上数据可估算出地球表面大气在标准状况下的体积为\xzanswer{B}

\fourchoices[answer=B]
{$4\times10^{16}\Umc$}
{$4\times10^{18}\Umc$}
{$4\times10^{20}\Umc$}
{$4\times10^{22}\Umc$}

%% answer: B
%% memo:
\memoanswer{
大气压是由大气重量产生的。大气压强 $p=\frac{mg}{S}=\frac{mg}{4\pi R^2}$，带入数据可得地球表面大气质量 $m=5.2\times10^{18}\Ukg$。标准状态下 $1\Umol$ 气体的体积为 $v=22.4\times10^{-3}\Umc$，故地球表面大气体积为 $V=\frac{m}{m_0}v=\frac{5.2\times10^{18}}{29\times10^{-3}}\times22.4\times10^{-3}\Umc=4\times10^{18}\Umc$，B 对。
}

\item
%% number: 20
%% paperName: 2008年全国理综Ⅰ第 20 题 6 分
%% typeId: 1
%% type: 单选题
%% chapter: 电磁感应
%% point: 电磁感应中图象
%% method: 无
%% score: 6
%% degree: 450
%% duplicateId: 0
%% body:
矩形导线框 abcd 固定在匀强磁场中，磁感线的方向与导线框所在平面垂直，规定磁场的正方向垂直纸面向里，磁感应强度 $B$ 随时间变化的规律如图所示。若规定顺时针方向为感应电流 $I$ 的正方向，下列各图中正确的是\xzanswer{D}

\onepicture[w=7.5cm]{figs/20a.png}

\fourchoices[answer=D, ispicture=true, h=1.8cm]
{figs/20b.png}
{figs/20c.png}
{figs/20d.png}
{figs/20e.png}

%% answer: D
%% memo:
\memoanswer{
$0\sim1\Us$ 内 $B$ 垂直纸面向里均匀增大，则由楞次定律及法拉第电磁感应定律可得线圈中产生恒定的感应电流，方向为逆时针方向，排除 A、C 选项；$2\Us\sim3\Us$ 内，$B$ 垂直纸面向外均匀增大，同理可得线圈中产生的感应电流方向为顺时针方向，排除 B 选项，D 正确。
}

\item
%% number: 21
%% paperName: 2008年全国理综Ⅰ第 21 题 6 分
%% typeId: 1
%% type: 单选题
%% chapter: 光学
%% point: 光的折射
%% method: 无
%% score: 6
%% degree: 450
%% duplicateId: 0
%% body:
一束由红、蓝两单色光组成的光线从一平板玻璃砖的上表面以入射角 $\theta$ 射入，穿过玻璃砖自下表面射出。已知该玻璃对红光的折射率为 1.5，设红光与蓝光穿过玻璃砖所用的时间分别为 $t_{1}$ 和 $t_{2}$，则在 $\theta$ 从 $0^{\circ}$ 逐渐增大至 $90^{\circ}$ 的过程中\xzanswer{B}

\fourchoices[answer=B]
{$t_{1}$ 始终大于 $t_{2}$}
{$t_{1}$ 始终小于 $t_{2}$}
{$t_{1}$ 先大于后小于 $t_{2}$}
{$t_{1}$ 先小于后大于 $t_{2}$}

%% answer: B
%% memo:
\memoanswer{
设折射角为 $\alpha$，玻璃砖的厚度为 $h$，由折射定律 $n=\frac{\sin\theta}{\sin\alpha}$，且 $n=\frac{c}{v}$，在玻璃砖中的时间为 $t=\frac{h}{v\cos\alpha}$，联立解得 $t_2\propto\frac{n^4}{n^2-\sin^2\theta}$，红光频率较小，$\theta$ 为零时，$t_1<t_2$，$\theta$ 为 $90^{\circ}$ 时，趋近渐近线，初步判定该函数为单调函数，通过带入 $\theta$ 为其它特殊值，仍然有 $t_1<t_2$，故 B 对。
}

\item
%% number: 22
%% paperName: 2008年全国理综Ⅰ第 22 题 12 分
%% typeId: 4
%% type: 实验题
%% chapter: 电路
%% point: 电路实验
%% method: 无
%% score: 12
%% degree: 450
%% duplicateId: 0
%% body:
\begin{enumerate}
	\item
	Ⅰ．（6 分）如图所示，两个质量各为 $m_{1}$ 和 $m_{2}$ 的小物块 A 和 B，分别系在一条跨过定滑轮的软绳两端，已知 $m_{1}>m_{2}$，现要利用此装置验证机械能守恒定律。

	（1）若选定物块 A 从静止开始下落的过程中进行测量，则需要测量的物理量有 \tkanswer[2.5cm]{}。

	①物块的质量 $m_{1}$、$m_{2}$；

	②物块 A 下落的距离及下落这段距离所用的时间；

	③物块 B 下落的距离及下落这段距离所用的时间；

	④绳子的长度。

	（2）为提高实验结果的准确程度，某小组同学对此实验提出以下建议：

	①绳的质量要轻；

	②在“轻质绳”的前提下，绳子越长越好；

	③尽量保证物块只沿竖直方向运动，不要摇晃；

	④两个物块的质量之差要尽可能小。

	以上建议中确实对提高准确程度有作用的是 \tkanswer[2.5cm]{}。

	（3）写出一条上面没有提到的对提高实验结果准确程度有益的建议：\tkanswer[4cm]{}。

	\onepicture[h=5.0cm, align=right]{figs/22a.png}

	\item
	Ⅱ．（12 分）一直流电压表 V，量程为 $1\UV$，内阻为 $1000\UO$。现将一阻值在 $5000\sim7000\UO$ 之间的固定电阻 $R_{1}$ 与此电压表串联，以扩大电压表的量程。为求得扩大后量程的准确值，再给定一直流电源（电动势 $E$ 为 $6\sim7\UV$，内阻可忽略不计），一阻值 $R_{2}=2000\UO$ 的固定电阻，两个单刀开关 $S_{1}$、$S_{2}$ 及若干导线。

	（1）为达到上述目的，将答题卡上对应的图连成一个完整的实验电路图；

	（2）连线完成以后，当 $S_{1}$ 与 $S_{2}$ 均闭合时，电压表的示数为 $0.90\UV$；当 $S_{1}$ 闭合，$S_{2}$ 断开时，电压表的示数为 $0.70\UV$，由此可以计算出改装后电压表的量程为 \tkanswer[1.5cm]{} $\UV$，电源电动势为 \tkanswer[1.5cm]{} $\UV$。

	\onepicture[h=5.5cm, align=right]{figs/22b.png}
\end{enumerate}

%% answer:
\jdanswer{
\begin{enumerate}
	\item
	Ⅰ．（1）①②或①③；（2）①③；（3）对同一高度进行多次测量取平均值；或选取受力后相对伸长量尽量小的绳。

	\item
	Ⅱ．将待测电压表与标准电阻串联后与电源连接即可，电路图如图所示。改装后电压表的量程为 $7\UV$，电源电动势为 $6.3\UV$。

	\onepicture[w=6.0cm]{figs/22_an.png}
\end{enumerate}
}

%% memo:
\memoanswer{
Ⅰ．（1）通过连结在一起的 A、B 两物体验证机械能守恒定律，即验证系统的势能变化与动能变化是否相等。A、B 连结在一起，A 下降的距离一定等于 B 上升的距离；A、B 的速度大小总是相等的，故不需要测量绳子的长度和 B 上升的距离及时间。（2）如果绳子质量不能忽略，则 A、B 组成的系统势能将有一部分转化为绳子的动能，从而为验证机械能守恒定律带来误差；若物块摇摆，则两物体的速度有差别，为计算系统的动能带来误差；绳子长度和两个物块质量差应适当。（3）对同一高度进行多次测量取平均值；或选取受力后相对伸长量尽量小的绳。

Ⅱ．将待测电压表与标准电阻串联后与电源连接即可。设电源电动势为 $E$，则由闭合电路欧姆定律，当两开关都闭合时 $R_{2}$ 被短路，有 $U_{1}=E$；当 $S_{1}$ 闭合、$S_{2}$ 断开时，$E=U_{2}+\frac{U_{2}}{R_{0}}(R_{1}+R_{2})$；解两式得 $R_{1}=6000\UO$，$E=6.3\UV$；根据串联分压原理，可得改装后电压表的量程为 $7\UV$。
}

\item
%% number: 23
%% paperName: 2008年全国理综Ⅰ第 23 题 14 分
%% typeId: 6
%% type: 计算题
%% chapter: 直线运动
%% point: 多段匀变速直线运动
%% method: 无
%% score: 14
%% degree: 650
%% duplicateId: 0
%% body:
已知 O、A、B、C 为同一直线上的四点，AB 间的距离为 $l_{1}$，BC 间的距离为 $l_{2}$，一物体自 O 点由静止出发，沿此直线做匀加速运动，依次经过 A、B、C 三点，已知物体通过 AB 段与 BC 段所用的时间相等。求 O 与 A 的距离。

%% answer:
\jdanswer{$l=\frac{(3l_{1}-l_{2})^{2}}{8(l_{2}-l_{1})}$}

%% memo:
\memoanswer{
设物体的加速度为 $a$，到达 A 点的速度为 $v_{0}$，通过 AB 段和 BC 段所用的时间为 $t$，则有 $l_{1}=v_{0}t+\frac{1}{2}at^{2}$ ①，$l_{1}+l_{2}=2v_{0}t+2at^{2}$ ②。联立①②式得 $l_{2}-l_{1}=at^{2}$ ③，$3l_{1}-l_{2}=2v_{0}t$ ④。设 O 与 A 的距离为 $l$，则有 $l=\frac{v_{0}^{2}}{2a}$ ⑤。联立③④⑤式得 $l=\frac{(3l_{1}-l_{2})^{2}}{8(l_{2}-l_{1})}$。
}

\item
%% number: 24
%% paperName: 2008年全国理综Ⅰ第 24 题 18 分
%% typeId: 6
%% type: 计算题
%% chapter: 运动和力
%% point: 动量守恒定律
%% method: 无
%% score: 18
%% degree: 450
%% duplicateId: 0
%% body:
图中滑块和小球的质量均为 $m$，滑块可在水平放置的光滑固定导轨上自由滑动，小球与滑块上的悬点 O 由一不可伸长的轻绳相连，轻绳长为 $l$。开始时，轻绳处于水平拉直状态，小球和滑块均静止。现将小球由静止释放，当小球到达最低点时，滑块刚好被一表面涂有粘性物质的固定挡板粘住，在极短的时间内速度减为零，小球继续向左摆动，当轻绳与竖直方向的夹角 $\theta=60^{\circ}$ 时小球达到最高点。求：

\begin{enumerate}
	\item
	从滑块与挡板接触到速度刚好变为零的过程中，挡板阻力对滑块的冲量；
	\item
	小球从释放到第一次到达最低点的过程中，绳的拉力对小球做功的大小。
\end{enumerate}

\onepicture[w=6.0cm, align=right]{figs/24.png}

%% answer:
\jdanswer{
\begin{enumerate}
	\item
	$I=m\sqrt{gl}$，方向向左
	\item
	$W=-\frac{1}{2}mgl$
\end{enumerate}
}

%% memo:
\memoanswer{
（1）对系统，设小球在最低点时速度大小为 $v_{1}$，此时滑块的速度大小为 $v_{2}$，滑块与挡板接触前由系统的机械能守恒定律：$mgl=\frac{1}{2}mv_{1}^{2}+\frac{1}{2}mv_{2}^{2}$ ①；由系统的水平方向动量守恒定律：$mv_{1}=mv_{2}$ ②；对滑块与挡板接触到速度刚好变为零的过程中，挡板阻力对滑块的冲量为 $I=mv_{2}$ ③。联立①②③解得 $I=m\sqrt{gl}$，方向向左 ④。

（2）小球释放到第一次到达最低点的过程中，设绳的拉力对小球做功的大小为 $W$，对小球由动能定理：$mgl+W=\frac{1}{2}mv_{1}^{2}$ ⑤。联立①②⑤解得 $W=-\frac{1}{2}mgl$，即绳的拉力对小球做负功，大小为 $\frac{1}{2}mgl$。
}

\item
%% number: 25
%% paperName: 2008年全国理综Ⅰ第 25 题 22 分
%% typeId: 6
%% type: 计算题
%% chapter: 磁场
%% point: 带电粒子在磁场中的运动
%% method: 无
%% score: 22
%% degree: 450
%% duplicateId: 0
%% body:
如图所示，在坐标系 $xOy$ 中，过原点的直线 OC 与 $x$ 轴正向的夹角 $\varphi=120^{\circ}$，在 OC 右侧有一匀强电场；在第二、三象限内有一匀强磁场，其上边界与电场边界重叠、右边界为 $y$ 轴、左边界为图中平行于 $y$ 轴的虚线，磁场的磁感应强度大小为 $B$，方向垂直纸面向里。一带正电荷 $q$、质量为 $m$ 的粒子以某一速度自磁场左边界上的 A 点射入磁场区域，并从 O 点射出，粒子射出磁场的速度方向与 $x$ 轴的夹角 $\theta=30^{\circ}$、大小为 $v$，粒子在磁场中的运动轨迹为纸面内的一段圆弧，且弧的半径为磁场左右边界间距的两倍。粒子进入电场后，在电场力的作用下又由 O 点返回磁场区域，经过一段时间后再次离开磁场。已知粒子从 A 点射入到第二次离开磁场所用的时间恰好等于粒子在磁场中做圆周运动的周期。忽略重力的影响。求：

\begin{enumerate}
	\item
	粒子经过 A 点时速度的方向和 A 点到 $x$ 轴的距离；
	\item
	匀强电场的大小和方向；
	\item
	粒子从第二次离开磁场到再次进入电场时所用的时间。
\end{enumerate}

\onepicture[w=5.5cm, align=right]{figs/25.png}

%% answer:
\jdanswer{
\begin{enumerate}
	\item
	速度方向沿 $x$ 轴正方向，A 点到 $x$ 轴的距离为 $\frac{(2-\sqrt{3})mv}{2Bq}$
	\item
	$E=\frac{12Bv}{7\pi}$，方向垂直于 OC 向左
	\item
	$t_{2}=\frac{\sqrt{3}m}{Bq}$
\end{enumerate}
}

%% memo:
\memoanswer{
（1）从 A 点进入磁场后从 O 点离开磁场的过程是匀速圆周运动，画出粒子运动的轨迹图，依题意由几何关系可得圆弧的圆心正好是两条虚线的交点，故经过 A 点的速度方向为 $x$ 轴正方向。设圆周的半径为 $R$，有 $\angle OO_{1}A=30^{\circ}$ ①；根据向心力公式 $Bqv=m\frac{v^{2}}{R}$ ②；A 点到 $x$ 轴的距离 $x=R-R\cos30^{\circ}$ ③。联立①②③解得 $x=\frac{(2-\sqrt{3})mv}{2Bq}$。

\onepicture[w=4.2cm]{figs/25_an1.png}

（2）粒子能从 O 点进入电场且能由 O 点返回，对正电荷，说明电场的方向垂直于 OC 向左，设电场强度大小为 $E$，电场中的时间为 $t_{1}$，由动量定理 $Eqt_{1}=2mv$ ④。粒子从 A 点射入到第二次离开磁场所用的时间恰好等于粒子在磁场中做圆周运动的周期 $T$，由 $T=\frac{2\pi R}{v}$ ⑤。从 O 点返回磁场后的轨迹如图，圆心角为 $120^{\circ}$，故 $T=t_{1}+\frac{1}{12}T+\frac{1}{3}T$ ⑥。联立②④⑤⑥解得 $E=\frac{12Bv}{7\pi}$ ⑦，方向垂直于 OC 向左。

\onepicture[w=3.2cm]{figs/25_an2.png}

（3）第二次离开磁场后到再进入电场，如图轨迹，则 $DF=OD=2R\cos30^{\circ}$ ⑧，时间 $t_{2}=\frac{OD}{v}=\frac{\sqrt{3}m}{Bq}$。

\onepicture[w=4.0cm]{figs/25_an3.png}
}

\end{enumerate}

\end{document}
